\documentclass[11pt]{article}
\usepackage{amsmath,amssymb,amsthm,mathtools}
\usepackage[a4paper,margin=1in]{geometry}
\usepackage{bm}
\usepackage{xcolor}
\usepackage{hyperref}

\newtheorem{theorem}{Theorem}
\newtheorem{remark}[theorem]{Remark}

\title{An HDG-type Discretization of 1D Embedded Blood Flow in Complex Networks with Riemann-Coupled Junctions}
\author{}
\date{}

\newcommand{\R}{\mathbb{R}}
\newcommand{\dt}{\Delta t}
\newcommand{\Th}{\mathcal{T}_h}
\newcommand{\Eh}{\mathcal{E}_h}
\newcommand{\Eho}{\mathcal{E}_h^o}
\newcommand{\Ehb}{\mathcal{E}_h^{\partial}}
\newcommand{\Ehj}{\mathcal{E}_h^{J}}
\newcommand{\norm}[1]{\left\lVert #1\right\rVert}
\newcommand{\wh}{\mathbf{w}_h}
\newcommand{\whl}{\mathbf{w}^L_h}
\newcommand{\whr}{\mathbf{w}^R_h}
\newcommand{\who}{\widehat{\mathbf{w}}_h}
\newcommand{\Hn}{\mathbf{H}_n}
\newcommand{\Hnl}{\mathbf{H}_n^{L}}
\newcommand{\Hnr}{\mathbf{H}_n^{R}}

\begin{document}
\maketitle
\section{Introduction}
\paragraph{Motivation for HDG.}
\begin{itemize}
\item 
Discontinuous Galerkin (DG) methods are popular for hyperbolic conservation laws because they handle shocks and complex geometry well and parallelize nicely. However, in a standard DG formulation, the numerical flux on each element depends on its neighboring element states, leading to a globally coupled system involving all element degrees of freedom. Solving that requires inverting one enormous global system. 

\item Hybridizable discontinuous Galerkin (HDG) methods, fix this by introducing extra "trace" unknowns living only on the mesh skeleton (namely the trace of the unknown on the mesh faces). Once those trace values are known, each element can be solved completely independently through local element solver \cite{MR3345205}.
\end{itemize}
\paragraph{Contribution.}
Hybridizable discontinuous Galerkin (HDG) methods have
been extensively studied over the last two decades, and
several HDG formulations incorporating approximate
Riemann solvers such as the Harten--Lax--van Leer (HLL)
flux have been proposed for hyperbolic systems,
particularly in the context of compressible fluid
dynamics \cite{peraire2010hybridizable,vila2021hybridisable}.
However, to the best of the authors' knowledge, there is
very limited work on HDG--HLL formulations for
one-dimensional blood flow models, especially in vascular
networks containing bifurcations and junctions.

The main challenge arises from the treatment of
junctions. While standard HDG formulations naturally
handle element interfaces through single-valued trace
variables, vascular junctions require the simultaneous
enforcement of mass conservation and pressure continuity
across multiple connected vessels. The present work
develops an HDG--HLL formulation for one-dimensional
blood flow networks that incorporates a dedicated
junction treatment within the hybridized framework

\section{Governing System}
We consider the one-dimensional blood flow equations parameterized along a centerline tangent vector $\bm b(\bm x)$. The compact vector form of the conservation system is given by:
\begin{equation}
  \partial_t \mathbf{w} + \nabla \cdot \left(\bm b(\bm{x})  \mathbf{F}(\mathbf{w})\right) = \mathbf{S}(\mathbf{w}), \qquad
  \mathbf{w} = \begin{pmatrix} A \\ U \end{pmatrix}, \qquad
  \mathbf{S}(\mathbf{w}) = \begin{pmatrix} 0 \\ - \frac{f U}{A} \end{pmatrix},
\end{equation}
where $A(x,t)$ is the cross-sectional area, $U(x,t)$ is the mean velocity, $\rho$ is the blood density, and
$f = 2(\zeta+2)\,\mu\,\pi$ is the viscous friction coefficient
with dynamic viscosity $\mu$ and velocity-profile parameter
$\zeta$. 
The non-linear physical flux vector $\mathbf{F}(\mathbf{w})$ is:
\begin{equation}
  \mathbf{F}(\mathbf{w}) = \begin{pmatrix} F_A(\mathbf{w}) \\ F_U(\mathbf{w}) \end{pmatrix} = \begin{pmatrix} A U \\ \frac{U^2}{2} + \frac{1}{\rho} P(A) \end{pmatrix}.
\end{equation}

Note: In the code we are taking for now $\bm b (\bm x) \equiv \bm b.$

\subsection{Tube Law}\label{sec:tube_law}
The arterial wall mechanics are governed by the tube law~\nocite{matthys2007pulse}
\begin{align*}
P(A,x) - P_0 = P_d + \frac{\beta}{A_d(x)}
\left[ (A^m - A_d(x)^m) \right], \text{ with }  \beta = 4/3 \sqrt{\pi} E h_w(x),
\end{align*}
where $P_0 $ is the external pressure, $P_d$ is the diastolic pressure,
$E$ is the wall Young's modulus,
$h_w(x)$ is the wall thickness, and
$A_d(x)$ is the \emph{local diastolic area}, tube law constant $m =0.5$.

\paragraph{Tapered vessels.}
For vessels whose cross-section varies along the centreline
(as in the 37 arteries network and ADAN56 model~\cite{alastruey2012arterial, boileau2015benchmark}),
each vessel segment is characterised by a proximal (inlet) radius
$r_{\mathrm{in}}$ and a distal (outlet) radius $r_{\mathrm{out}}$.
The diastolic radius at arc-length coordinate $s$ is
\begin{equation}
  r_d(s) = r_{\mathrm{in}}
  + \frac{s - s_{\min}}{s_{\max} - s_{\min}}
  \bigl(r_{\mathrm{out}} - r_{\mathrm{in}}\bigr),
  \qquad
  A_d(s) = \pi\,r_d(s)^2,
  \label{eq:taper}
\end{equation}
where $s_{\min}$ and $s_{\max}$ are the arc-length bounds of the
vessel.
For a uniform vessel, $r_{\mathrm{in}} = r_{\mathrm{out}} = r_d$
and $A_d$ reduces to the constant $\pi r_d^2$.

The pressure derivative, wave speed, and wave-speed derivative are:
\begin{align}
  \frac{\partial P}{\partial A}\bigg|_x
  &= \frac{\beta(x)}{2\,A_d(x)\,\sqrt{A}},
  \label{eq:dPdA}
  \\[4pt]
  c(A,x)
  &= \sqrt{\frac{A}{\rho}\,\frac{\partial P}{\partial A}}
  = \sqrt{\frac{\beta(x)}{2\,\rho\,A_d(x)}}\;A^{1/4},
  \label{eq:wave_speed}
  \\[4pt]
  \frac{\partial c}{\partial A}
  &= \frac{c(A,x)}{4\,A}.
  \label{eq:dc_dA}
\end{align}

\begin{remark}[\textbf{As applied in the code}]
Mainly in the literature we are given $r_{in}$ and $r_{out}$ based on this we need to find out the linear varying $r_d$. In the discrete scheme the diastolic area $A_d$ appears in
three distinct places:
\begin{enumerate}
\item \textbf{Volume integrals:} $A_d$ is evaluated at the cell
center.
\item \textbf{Interior state at faces:}
The interior states are evaluated using the same
cell-wise diastolic area employed in the volume
integrals. Thus, for an interior face shared by
cells $K^-$ and $K^+$, we use
\[
A_d^- = A_d(x_{K^-}^{\mathrm{center}}),
\qquad
A_d^+ = A_d(x_{K^+}^{\mathrm{center}}),
\]
so that
\[
P^- = P(A^-,A_d^-),
\qquad
P^+ = P(A^+,A_d^+).
\]
This is consistent with the piecewise-constant
approximation of $A_d$ within each cell.
\item \textbf{Trace state at faces:}
Since the trace unknown $\widehat{A}$ is defined on the
face, the corresponding diastolic area is evaluated at
the shared face vertex $x_{\mathrm{face}}$. Let
\[
s_{\mathrm{face}}
=
x_{\mathrm{face}}\cdot \widehat{\mathbf d},
\]
where $\widehat{\mathbf d}$ is the vessel direction.
Then
\[
A_d^{\mathrm{face}}
=
A_d(s_{\mathrm{face}}).
\]
\end{enumerate}
\end{remark}

\subsection{Characteristic Structure}

For the physical flux
\[
  \mathbf{F}(\mathbf{w}) =
  \begin{pmatrix}
    A\,U \\[1mm]
    \dfrac{U^2}{2} + \dfrac{1}{\rho}P(A)
  \end{pmatrix},
\]
we have jacobian
\[
  \mathbf{H}(\mathbf{w})
  = \frac{\partial\mathbf{F}}{\partial\mathbf{w}}
  =
  \begin{pmatrix}
    \dfrac{\partial F_A}{\partial A} &
    \dfrac{\partial F_A}{\partial U} \\[2mm]
    \dfrac{\partial F_U}{\partial A} &
    \dfrac{\partial F_U}{\partial U}
  \end{pmatrix}
  =
  \begin{pmatrix}
    U & A \\[2mm]
    \dfrac{1}{\rho}\dfrac{dP}{dA}(A) & U
  \end{pmatrix}
  =
  \begin{pmatrix}
    U & A \\[2mm]
    \dfrac{c^2}{A} & U
  \end{pmatrix},
\]
where $c$ is the local wave speed.

The Jacobian of the physical flux  has eigenvalues
$\lambda_{f,b} = U \pm c$, defining forward and backward
characteristic paths.
Under physiological conditions the flow is subcritical
($|U| < c$), so information propagates in both directions.
The forward and backward Riemann invariants are
\begin{equation}
  W_f = U + 4(c - c_0),
  \qquad
  W_b = U - 4(c - c_0),
  \label{eq:riemann_inv}
\end{equation}
where $c_0 = c(A_d, x)$ is the initial wave speed at start of each vessel. Throughout the remainder of this work, the forward and backward characteristic variables are equivalently denoted by
\[
W_1 := W_f,
\qquad
W_2 := W_b,
\]
with indices $1$ and $2$ corresponding to the forward and backward characteristic fields, respectively.
\section{Hybridizable Discontinuous Galerkin Framework}
Let $\Th$ be the collection of non-overlapping line elements $K$, and let $\Eh$ be the set of all boundary endpoints (nodes)~\cite{peraire2010hybridizable}. We partition $\Eh$ into three mutually disjoint subsets:
\[
  \Eh = \Eho \cup \Ehb \cup \Ehj,
\]
representing standard interior nodes ($\Eho$), exterior boundary terminals ($\Ehb$), and branching junctions ($\Ehj$). 

\paragraph{Local and trace variables.}
We define the local element state space $\mathbf{W}_h^k$ and the trace space $\mathbf{M}_h^k$ as:
\begin{align*}
  \mathbf{W}_h^k &= \left\{ \mathbf{v} \in (L^2(\Th))^2 : \mathbf{v}|_K \in (\mathcal{P}^k(K))^2, \, \forall K \in \Th \right\}, \\
  \mathbf{M}_h^k &= \left\{ \bm{\mu} \in (L^2(\Eh))^2 : \bm{\mu}|_F \in (\mathcal{P}^k(F))^2, \, \forall F \in \Eh \right\}.
\end{align*}
the HDG trace variable $\widehat{\mathbf{w}}_h = (\widehat{A}_h, \widehat{U}_h)^\top \in \mathbf{M}_h^k$ is defined over $\Eh$.

\paragraph{The HDG Weak Formulation.}
Inside each element $K \in \Th$, the interior state $\mathbf{w}_h \in \mathbf{W}_h^k$ is recovered locally by the trace unknown on $\partial K$. These trace unknowns are determined globally through the interface, boundary, and junction coupling conditions. The local cell equation reads: 
\begin{equation}
  \left(\partial_t \mathbf{w}_h, \mathbf{v}\right)_K - \left(\mathbf{F}(\mathbf{w}_h), \bm b \cdot \nabla \mathbf{v}\right)_K + \left\langle \widehat{\mathbf{F}}_h \cdot \bm n, \mathbf{v} \right\rangle_{\partial K} = \left(\mathbf{S}(\mathbf{w}_h), \mathbf{v}\right)_K, \quad \forall \mathbf{v} \in (\mathcal{P}^k(K))^2,
  \label{eq:hdg_local}
\end{equation}
where $\bm n$ is the outward unit normal vector pointing out of element $K$. 

\subsection{Numerical Trace Flux}
The numerical trace flux $\widehat{\mathbf F}_h\cdot\bm n$ on $\partial K$ is a
\emph{two-state interface flux}. On any face it receives the interior element
state $\mathbf w_h=(A_h,U_h)^\top$ (the \emph{self} state, evaluated on the face
from inside $K$) and the single-valued trace $\widehat{\mathbf w}_h=(\widehat
A_h,\widehat U_h)^\top$ (the \emph{exterior} state living on the face), and
returns a single normal flux,
\begin{equation}
  \widehat{\mathbf F}_h\cdot\bm n
  \;=\;
  \mathcal F_{\mathrm{num}}\bigl(\,b\!\cdot\!\bm n,\;\mathbf w_h,\;
  \widehat{\mathbf w}_h\,\bigr).
  \label{eq:numflux_two_state}
\end{equation}
Here, we write $\mathcal F_{\mathrm{num}}$ for every standard numerical flux choice like Lax--Friedrichs, Rusanov, Roe, or HLL as a unified form which satisfies the HDG form like:
\begin{equation}
  \widehat{\mathbf F}_h\cdot\bm n
  \;=\;
  \underbrace{\mathbf F(\widehat{\mathbf w}_h)\cdot\bm n}_{\text{consistent part}}
  \;+\;
  \underbrace{\bm\Sigma(\mathbf w_h,\widehat{\mathbf w}_h)\,
  \bigl(\mathbf w_h-\widehat{\mathbf w}_h\bigr)}_{\text{stabilization / dissipation}},
  \label{eq:hdg_template}
\end{equation}
where the physical flux is anchored at the trace and the matrix
$\bm\Sigma$ penalizes the interior--trace mismatch and controls the amount of numerical dissipation and ensures stability of hybridized formulation.

The simplest closure is the scalar Local Lax--Friedrichs (LLF) penalty
$\bm\Sigma_{\mathrm{LF}}=\alpha_{\mathrm{LF}}\,\mathbf I$ with
$$\alpha_{\mathrm{LF}}=\max\bigl(|U_h\pm c_h|,\,|\widehat U_h\pm\widehat
c_h|\bigr),$$ which spreads identical dissipation over both characteristic
directions. 

Unlike scalar Lax–Friedrichs stabilization, the HLL flux uses separate estimates of the left- and right-running wave speeds. This allows the numerical flux to distinguish between the two characteristic directions and provide wave-dependent dissipation.

\paragraph{From the HDG stabilisation to the HLL flux.}
In the HDG context the self/interior state is the
left state and the trace is the right state,
\begin{equation}
  \mathbf w_L \equiv \mathbf w_h = (A_h,U_h)^\top,
  \qquad
  \mathbf w_R \equiv \widehat{\mathbf w}_h = (\widehat A_h,\widehat U_h)^\top .
\end{equation}
The orientation-projected physical fluxes are
\begin{equation}
  \mathbf F_\star\cdot\bm n
  =
  (b\!\cdot\!\bm n)
  \begin{pmatrix}
  A_\star U_\star \\[2pt]
  \tfrac12 U_\star^2 + P(A_\star)/\rho
  \end{pmatrix},
  \qquad \star\in\{L,R\},
  \label{eq:proj_flux}
\end{equation}
with the wave speed from the tube law $c(A)$.
The signal speeds are estimated, as by the average states
\begin{equation}
  \bar U=\tfrac12(U_L+U_R),
  \quad
  \bar c=\tfrac12\bigl(c(A_L)+c(A_R)\bigr),
  \qquad
  s_L=\bar U-\bar c,
  \quad
  s_R=\bar U+\bar c .
  \label{eq:wave_speeds}
\end{equation}
The HLL numerical flux is then the three-branch expression \cite[Chapter~10]{toro2013riemann}
\begin{equation}
  \widehat{\mathbf F}_h\cdot\bm n
  =
  \begin{cases}
  \mathbf F_R\cdot\bm n, & s_R\le 0,\\[6pt]
  \dfrac{s_R\,\mathbf F_L\cdot\bm n - s_L\,\mathbf F_R\cdot\bm n
         + s_R s_L\,(\mathbf w_R-\mathbf w_L)}{s_R-s_L}, 
  & s_L<0<s_R, \\[6pt]
  \mathbf F_L\cdot\bm n, & s_L\ge 0,
  \end{cases}
  \label{eq:hll_flux}
\end{equation}
which is exactly the returned object of \texttt{hll\_flux} in the code: the supersonic
branches (for which $|U| >c$) pick the pure upwind flux, and the subcritical branch
$|U|<c$, the physiological regime, always selected for blood flow, is with the dissipative jump $s_R s_L(\mathbf w_R-\mathbf w_L)$
carried on the \emph{raw} states.

\paragraph{Equivalence with the HDG normal flux.}
The central branch of~\eqref{eq:hll_flux} is itself the numerical flux, and decompose it into a consistent part plus a \emph{scalar}
dissipation acting on the interior--trace jump. Splitting the
numerator,
\begin{equation}
  \widehat{\mathbf F}_h\cdot\bm n
  =
  \underbrace{\frac{s_R\,\mathbf F(\mathbf w_h)\cdot\bm n
            -s_L\,\mathbf F(\widehat{\mathbf w}_h)\cdot\bm n}{s_R-s_L}}
            _{\text{consistent (wave-weighted) flux}}
  \;+\; \underbrace{\bm\Sigma_{\mathrm{HLL}}\,
  \bigl(\mathbf w_h-\widehat{\mathbf w}_h\bigr)}_{\text{dissipation}}, 
  \label{eq:hll_consistent}
\end{equation}
where $$
  \bm\Sigma_{\mathrm{HLL}}=\delta_{\mathrm{HLL}}\,\mathbf I,
  \quad
  \delta_{\mathrm{HLL}}= -\frac{\,s_L\,s_R}{\,s_R-s_L\,}\;\ge 0 .$$
Here the stabilisation $\bm\Sigma_{\mathrm{HLL}}$ is the \emph{scalar} matrix
$\delta_{\mathrm{HLL}}\mathbf I$: in subcritical flow $s_L<0<s_R$, so
$\delta_{\mathrm{HLL}}=-s_Ls_R/(s_R-s_L)>0$. 

Equation~\eqref{eq:hll_consistent}is identical to
the central branch returned by \texttt{hll\_flux} (with $\mathbf
w_h\equiv(A_L,U_L)$ the interior/self state and $\widehat{\mathbf
w}_h\equiv(A_R,U_R)$ the trace). The implementation uses the mean-state estimate~\eqref{eq:wave_speeds}; the
Roe bounds $s_L= \bar{U} - \bar{c}$, $s_R=\bar{U} + \bar{c}$, with $\bar{U} = 0.5(U_L + U_R) $ and $\bar{c} = 0.5(c(A_L)+c(A_R))$.
 

\begin{remark}[HLL $\to$ Rusanov $\to$ LLF]
With a single symmetric speed $s_R=-s_L=\alpha$, the average
in~\eqref{eq:hll_consistent} becomes $\tfrac12\bigl(\mathbf F(\mathbf
w_h)+\mathbf F(\widehat{\mathbf w}_h)\bigr)\cdot\bm n$ and
$\delta_{\mathrm{HLL}}=\alpha/2$, recovering the Rusanov / local
Lax--Friedrichs flux
$\tfrac12\bigl(\mathbf F(\mathbf w_h)+\mathbf F(\widehat{\mathbf
w}_h)\bigr)\cdot\bm n+\tfrac{\alpha}{2}\bigl(\mathbf w_h-\widehat{\mathbf
w}_h\bigr)$. HLL is therefore the two-wave generalisation that weights the
forward ($U+c$) and backward ($U-c$) characteristics separately instead of bounding both by one $\alpha$.
\end{remark}

\section{Unified Transmission Condition}

The global problem is closed by requiring conservative transmission of the numerical trace flux at every interface: find
$\widehat{\mathbf w}_h\in\mathbf M_h^k$ such that
\begin{equation}
  \sum_{K\in\Th}
  \bigl\langle\widehat{\mathbf F}_h\cdot\bm n,\bm\mu\bigr\rangle
    _{\partial K\setminus(\Ehb\cup\Ehj)}
  +\sum_{F\in\Ehb}\bigl\langle\widehat{\mathbf B}_h,\bm\mu\bigr\rangle_F
  +\sum_{J\in\Ehj}\bigl\langle\widehat{\mathbf J}_h,\bm\mu\bigr\rangle_J
  =0,
  \qquad\forall\bm\mu\in\mathbf M_h^k .
  \label{eq:global_hdg}
\end{equation}

The first term represents conservation of the numerical HLL flux across
standard interior interfaces. The operators
$\widehat{\mathbf B}_h$ and $\widehat{\mathbf J}_h$
represent, respectively, the contributions arising from physical boundary
conditions and vascular junction coupling conditions. Their construction is
presented in the following sections.

A key feature of the formulation is that the same HLL numerical
flux~\eqref{eq:hll_flux} is employed both in the local cell problem and in
the global transmission condition. This provides a consistent coupling
between the element-interior solution and the globally defined trace
unknowns, while allowing boundary and junction physics to be incorporated
through the operators $\widehat{\mathbf B}_h$ and
$\widehat{\mathbf J}_h$.
\subsection{Interior faces (HLL conservation)}

For an interior face shared by $K_L,K_R$ of the \emph{same} vessel, with
$b\!\cdot\!\bm n_R=-\,b\!\cdot\!\bm n_L$ and a single-valued trace, the residual
is the sum of the two one-sided HLL fluxes, each taking that side's interior as
the self state and the shared trace as the exterior state:
\begin{equation}
  \mathcal F_{\mathrm{int}}
  =
  \widehat{\mathbf F}_h\bigl(b\!\cdot\!\bm n_L,\,\mathbf w_L,\,
  \widehat{\mathbf w}_h\bigr)
  +
  \widehat{\mathbf F}_h\bigl(b\!\cdot\!\bm n_R,\,\mathbf w_R,\,
  \widehat{\mathbf w}_h\bigr)
  =\mathbf 0 .
  \label{eq:interior_hll}
\end{equation}
These are two scalar equations (area and momentum rows) fixing the two trace
unknowns $(\widehat A_h,\widehat U_h)$. In the code this is the pair
\texttt{(FA\_L+FA\_R, FU\_L+FU\_R)} returned by two calls to the HLL flux with
$\mathtt{bn\_L}$ and $\mathtt{bn\_R}=-\mathtt{bn\_L}$, the trace passed as the
exterior state on both sides, and each face assembled once and given in \texttt{assemble\_trace\_interior\_eqautions()} for the residual case and the jacobian case is in \texttt{assemble\_jacobian\_trace\_interior\_eqautions()}.

\subsection{Physical boundaries (characteristic operator $\widehat{\mathbf
B}_h$)}

At a boundary face there is no neighbour, so the trace cannot be closed by flux
conservation. For subcritical flow exactly one characteristic leaves the domain
and one physical condition may be imposed; $\widehat{\mathbf B}_h$ is therefore
[one physical condition]~$+$~[one outgoing Riemann-invariant compatibility],
with $c_0=c(A_d)$ and
\begin{equation}
  W_1 = U + 4(c-c_0)\ \ (\text{forward}),
  \qquad
  W_2 = U - 4(c-c_0)\ \ (\text{backward}),
\end{equation}
in code this is given in \texttt{assemble\_trace\_boundary\_equations()} and for the jacobian block \texttt{assemble\_jacobian\_trace\_boundary\_equations()}.


\paragraph{Inflow ($\mathtt{bid}=0$).} A flow rate $Q_{\mathrm{in}}(t)$ is
prescribed; the outgoing backward invariant is matched to the interior:
\begin{equation}
  \widehat{\mathbf B}_h^{\mathrm{in}}
  =
  \begin{pmatrix}
  \widehat A_h\,\widehat U_h - Q_{\mathrm{in}}(t)\\[6pt]
  \bigl[\widehat U_h-4(\widehat c_h-c_0)\bigr]
  -\bigl[U_h-4(c_h-c_0)\bigr]
  \end{pmatrix}
  =\mathbf 0 .
  \label{eq:bc_inflow}
\end{equation}
The Jacobian blocks are:
\begin{equation}
D_{\mathrm{in}}
=
\frac{\partial\widehat{\mathbf B}_h^{\mathrm{in}}}
     {\partial\widehat{\mathbf{w}}_h}
=
\begin{bmatrix}
\widehat{U}_h & \widehat{A}_h
\\[8pt]
4\,\dfrac{dc}{d\widehat{A}_h} & 1
\end{bmatrix},
\qquad
C_{\mathrm{in}}
=
\frac{\partial\widehat{\mathbf B}_h^{\mathrm{in}}}
     {\partial\mathbf{w}_h}
=
\begin{bmatrix}
0 & 0
\\[8pt]
4\,\dfrac{dc}{dA_h} & -1
\end{bmatrix},
\label{eq:inflow_jacobian}
\end{equation}

\paragraph{RCR / single-resistance outlet ($\mathtt{bid}\neq0$).} Pressure
continuity with the Windkessel terminal, plus the outgoing forward invariant:
\begin{equation}
  \widehat{\mathbf B}_h^{\mathrm{rcr}}
  =
  \begin{pmatrix}
  P(\widehat A_h) - \bigl(R_1\,\widehat A_h\widehat U_h + P_c\bigr)\\[6pt]
  \bigl[\widehat U_h+4(\widehat c_h-c_0)\bigr]
  -\bigl[U_h+4(c_h-c_0)\bigr]
  \end{pmatrix}
  =\mathbf 0 ,
  \label{eq:bc_rcr}
\end{equation}
where the capacitor pressure $P_c$ is carried as a global unknown
$P_c=\mathbf y[\,\texttt{rcr\_pc\_dof}\,]$ (fully coupled into the DAE). 

The
single-resistance case $C=0$ is the special case $P(\widehat A_h)-(R_2\,\widehat
A_h\widehat U_h+P_{\mathrm{out}})$ in the first row.

The Jacobian with respect to the trace variables is

\begin{equation}
D_{\mathrm{rcr}}
=
\frac{\partial \widehat{\mathbf B}_h^{\mathrm{rcr}}}
     {\partial \widehat{\mathbf w}_h}
=
\begin{bmatrix}
\dfrac{dP}{d\widehat A_h}
-
R_1\widehat U_h
&
-
R_1\widehat A_h
\\[10pt]
4\,\dfrac{dc}{d\widehat A_h}
&
1
\end{bmatrix}.
\label{eq:Drcr}
\end{equation}
The Jacobian with respect to the interior state is

\begin{equation}
C_{\mathrm{rcr}}
=
\frac{\partial \widehat{\mathbf B}_h^{\mathrm{rcr}}}
     {\partial \mathbf w_h}
=
\begin{bmatrix}
0 & 0
\\[8pt]
-4\,\dfrac{dc}{dA_h}
&
-1
\end{bmatrix}.
\label{eq:Brcr}
\end{equation}
since $P_c$ is a global degree of freedom, there is an additional coupling block
\begin{equation}
E_{\mathrm{rcr}}
=
\frac{\partial \widehat{\mathbf B}_h^{\mathrm{rcr}}}
     {\partial P_c}
=
\begin{bmatrix}
-1
\\[4pt]
0
\end{bmatrix},
\label{eq:Ercr}
\end{equation}
\paragraph{Terminal pressure.}
The capacitor pressure is determined by the additional DAE equation which is given by this  $C \ \dot P_c = A\;U - (P_c - P_{out})/ R2$  in code we write the residual as 
\begin{equation}
R_{pc} = C \dot P_c  -
\widehat A_h\widehat U_h
+
\frac{P_c-P_{\mathrm{out}}}{R_2}
\label{eq:rcr_capacitor}
\end{equation}
given in \texttt{assemble\_rcr\_capacitor\_equations()}
whose Jacobian contributions are
\begin{equation}
\frac{\partial R_{pc}}
{\partial (\widehat A_h,\widehat U_h,P_c)} =
\begin{bmatrix}
-\widehat U_h &
-\widehat A_h &
\dfrac{1}{R_2}
\end{bmatrix},
\label{eq:rcr_cap_jac}
\end{equation}
and in code \texttt{assemble\_jacobian\_rcr\_capacitor\_equations()}
together with the mass contribution 
\begin{equation}
\frac{\partial R_{pc}}
{\partial \dot P_c} = C.
\label{eq:rcr_cap_mass}
\end{equation}
in code in \texttt{assemble\_jacobian()} with \\\texttt{jacobian\_matrix.add(pc\_dof, pc\_dof, alpha * rcr\_map.at(bid).C)}

\paragraph{Reflection outlet.} A reflection coefficient $R_t$ sets the incoming
invariant from the outgoing one:
\begin{equation}
  \widehat{\mathbf B}_h^{\mathrm{refl}}
  =
  \begin{pmatrix}
  \bigl[\widehat U_h+4(\widehat c_h-c_0)\bigr]-W_{1}^{\mathrm{int}}\\[6pt]
  \bigl[\widehat U_h-4(\widehat c_h-c_0)\bigr]-\bigl(-R_t\,W_{1}^{\mathrm{int}}\bigr)
  \end{pmatrix}
  =\mathbf 0,
  \qquad
  W_1^{\mathrm{int}}=U_h+4(c_h-c_0).
  \label{eq:bc_refl}
\end{equation}
In all three cases the \emph{cell} equation on the same face still uses the HLL
flux $\widehat{\mathbf F}_h(\text{interior},\text{trace})$; only the trace value
is pinned by $\widehat{\mathbf B}_h$ instead of by conservation.

\subsubsection{Junctions (operator $\widehat{\mathbf J}_h$)}
Consider a junction $J$ connecting $K\ge 2$ vessels. 
\begin{itemize}
    \item For $K=2$, we define a two-way junction. In the implementation, if the two neighbouring elements sharing a face have the same vessel ID, the interface is treated as an interior trace. Otherwise, if the neighbouring elements belong to different vessels, the interface is classified as a junction and the corresponding junction conditions are imposed, which are given below.
\end{itemize}
Let
\[
\widehat{\mathbf{w}}_i
=
\begin{pmatrix}
\widehat{A}_i \\ \widehat{U}_i
\end{pmatrix},
\qquad i = 1,\ldots,K,
\]
denote the hybrid trace state of vessel $i$ at the junction, with
orientation sign
\[
s_i =
\begin{cases}
+1, & \text{if vessel } i \text{ ends at the junction},\\
-1, & \text{if vessel } i \text{ starts at the junction}.
\end{cases}
\]

Unlike physical boundaries, a junction is an internal network interface.
% interface. No external characteristic information is imposed.
% Instead, the standard HDG flux~\eqref{eq:cell_flux} is retained
% at each junction face:
% \begin{equation}
% \widehat{\mathbf{J}}_{h,i}
% =
% \mathbf{F}(\widehat{\mathbf{w}}_i)\cdot\bm{n}_i
% +
% \bm{\Sigma}_i(\mathbf{w}_{h,i} - \widehat{\mathbf{w}}_i),
% \qquad i = 1,\ldots,K.
% \label{eq:junction_flux}
% \end{equation}

The junction trace states are determined by the following
$K$ constraints, assembled as the transmission residual:
\begin{equation}
\mathcal{F}_J
(\widehat{\mathbf{w}}_1,\ldots,\widehat{\mathbf{w}}_K)
=
\begin{pmatrix}
\displaystyle\sum_{i=1}^{K}
s_i\,\widehat{A}_i\,\widehat{U}_i
\\[10pt]
\mathcal{H}(\widehat{\mathbf{w}}_1)
- \mathcal{H}(\widehat{\mathbf{w}}_2)
\\[4pt]
\vdots
\\[4pt]
\mathcal{H}(\widehat{\mathbf{w}}_1)
- \mathcal{H}(\widehat{\mathbf{w}}_K)
\end{pmatrix}
= \mathbf{0},
\qquad
\mathcal{H}(\widehat{\mathbf{w}}_i)
=
\tfrac{1}{2}\widehat{U}_i^2 + \Psi(\widehat{A}_i),
\label{eq:junction_transmission}
\end{equation}
where $\Psi(A) = P(A)/\rho$. Row 1 enforces global mass
conservation; rows 2 to $K$ enforce continuity of total energy
head between all vessels and the reference vessel $i=1$.

Since $\mathcal{F}_J$ depends only on junction trace unknowns
and not on any interior element state, the element-to-trace
coupling block vanishes:
\begin{equation}
C_{J,K_i}
=
\frac{\partial\mathcal{F}_J}{\partial\mathbf{w}_{h,i}}
=
\mathbf{0},
\qquad i=1,\ldots,K.
\label{eq:Cjunc}
\end{equation}

The trace control block is the $K\times 2K$ matrix:
\begin{equation}
D_J
=
\frac{\partial\mathcal{F}_J}
     {\partial(\widehat{\mathbf{w}}_1,\ldots,\widehat{\mathbf{w}}_K)}
=
\begin{bmatrix}
s_1\widehat{U}_1 & s_1\widehat{A}_1 &
s_2\widehat{U}_2 & s_2\widehat{A}_2 &
\cdots &
s_K\widehat{U}_K & s_K\widehat{A}_K
\\[8pt]
\Psi'(\widehat{A}_1) & \widehat{U}_1 &
-\Psi'(\widehat{A}_2) & -\widehat{U}_2 &
\cdots & 0 & 0
\\[4pt]
\vdots & & & & \ddots & & \vdots
\\[4pt]
\Psi'(\widehat{A}_1) & \widehat{U}_1 &
0 & 0 & \cdots &
-\Psi'(\widehat{A}_K) & -\widehat{U}_K
\end{bmatrix},
\label{eq:DJ}
\end{equation}
where $\Psi'(\widehat{A}_i) = (1/\rho)\,dP/d\widehat{A}_i$.
% The system at the junction is closed by the $B$ blocks from the
% element equations, which supply the remaining $K$ equations
% through the HDG numerical flux~\eqref{eq:junction_flux}.
\begin{remark}[Single-valued vs.\ multi-valued traces]
On interior faces and at physical boundaries the trace is single-valued, as in
standard HDG. At vessel junctions
this is relaxed: the cross-sectional area is physically discontinuous between
branches, so each incident face carries its own trace and the branch traces are
coupled through conservation of mass, continuity of total pressure and Riemann Invariants.
\end{remark}

\section{Solvability of semi-discrete HDG solution}
In this section we prove the solvability of semi-discrete HDG schemes which are given like this :
For any $(\mathbf{w}_h, \widehat{\mathbf{w}}_h\in \mathbf{W}_h^k \times \mathbf{M}_h^k$ we have 
\begin{subequations}
\begin{align}
    \left(\partial_t \mathbf{w}_h, \mathbf{v}\right)_K - \left(\mathbf{F}(\mathbf{w}_h), \bm b \cdot \nabla \mathbf{v}\right)_K + \left\langle \widehat{\mathbf{F}}_h \cdot \bm n, \mathbf{v} \right\rangle_{\partial K} &= \left(\mathbf{S}(\mathbf{w}_h), \mathbf{v}\right)_K, \quad \forall \mathbf{v} \in (\mathcal{P}^k(K))^2,
  \label{eq:semidiscrete1} \\
  \sum_{K \in \Th}
  \left\langle
  \widehat{\mathbf{F}}_h \cdot \bm n,
  \bm{\mu}
  \right\rangle_{\partial K \setminus (\Ehb \cup \Ehj)}
  +
  \sum_{F \in \Ehb}
  \left\langle
  \widehat{\mathbf{B}}_h,
  \bm{\mu}
  \right\rangle_F
  +
  \sum_{J \in \Ehj}
  \left\langle
  \widehat{\mathbf{J}}_h,
  \bm{\mu}
  \right\rangle_J
  &=0,
  \qquad
  \forall \bm{\mu}\in \mathbf{M}_h^k.
\label{eq:semidiscrete2} 
\end{align}
\end{subequations}
    

    


% ======================================================================
%  Time integration with SUNDIALS IDA.
%  Assumes the RCR boundary block already defines, in the boundary section:
%    eq:bc_rcr   (B_h^rcr trace residual)
%    eq:Drcr     (D_rcr  = dB/dwhat)
%    eq:Brcr     (C_rcr  = dB/dw_h)
%    eq:Ercr     (E_rcr  = dB/dPc = [-1;0])
%    eq:rcr_capacitor (R_pc, contains the C \dot P_c term)
%    eq:rcr_cap_jac   (dR_pc/d(Ahat,Uhat,Pc) = [-Uhat,-Ahat,1/R2])
%    eq:rcr_cap_mass  (dR_pc/d\dot P_c = C)
%  and the cell residual eq:RK, cell weak form eq:hdg_local.
% ======================================================================

\section{Time Integration with SUNDIALS IDA}

The differential--algebraic system is integrated with the \textsc{SUNDIALS} IDA
solver\footnote{\url{https://dealii.org/current/doxygen/deal.II/classSUNDIALS_1_1IDA.html}},
which advances a system of the form
\begin{equation}
  \mathbf F(t,\mathbf y,\dot{\mathbf y})=\mathbf 0
  \label{eq:ida_form}
\end{equation}
by variable-order, variable-step BDF (backward-differentiation formula)
methods, solving one nonlinear system per step by Newton's method.

\paragraph{Global unknown vector.}
The unknowns are grouped into a cell block, a trace block, and a capacitor
block,
\begin{equation}
  \mathbf y
  =
  \bigl(\mathbf y_{\mathrm{cell}},\;
        \mathbf y_{\mathrm{trace}},\;
        \mathbf y_{\mathrm{pc}}\bigr),
  \label{eq:ida_y}
\end{equation}
\begin{equation}
  \mathbf y_{\mathrm{pc}}
  =\bigl(P_{c,1},\dots,P_{c,N_{\mathrm{rcr}}}\bigr),
  \label{eq:ida_ypc}
\end{equation}
where $\mathbf y_{\mathrm{pc}}$ carries one capacitor pressure $P_{c,i}$ for
each \emph{capacitive} ($C_i>0$) RCR terminal; purely resistive ($C_i=0$)
outlets impose pressure algebraically inside $\mathbf R_{\mathrm{trace}}$ and
contribute no unknown.

\paragraph{Three-block residual.}
Stacking the cell weak form~\eqref{eq:hdg_local}, the transmission conditions,
and the Windkessel residual~\eqref{eq:rcr_capacitor} gives
\begin{equation}
  \mathbf F(t,\mathbf y,\dot{\mathbf y})
  =
  \begin{pmatrix}
  M_K\,\dot{\mathbf y}_{\mathrm{cell}}-\mathbf R_{\mathrm{cell}}(t,\mathbf y)\\[6pt]
  \mathbf R_{\mathrm{trace}}(t,\mathbf y)\\[6pt]
  \mathbf R_{\mathrm{pc}}(t,\mathbf y,\dot{\mathbf y}_{\mathrm{pc}})
  \end{pmatrix}
  =\mathbf 0 ,
  \label{eq:ida_residual}
\end{equation}
where $\mathbf R_{\mathrm{cell}}$ is the cell residual~\eqref{eq:RK} (volume
flux, face flux via the trace, and source), $\mathbf R_{\mathrm{trace}}$
collects the interior, inflow, single-resistance, RCR and junction transmission
residuals, and $\mathbf R_{\mathrm{pc}}$ is the capacitor
residual~\eqref{eq:rcr_capacitor}, which already contains the time derivative
$C_i\,\dot P_{c,i}$. The time derivative thus enters only through the cell mass
matrix $M_K$ and the capacitor block; the trace rows are purely algebraic.

\paragraph{Differential / algebraic split.}
The rows carrying a time derivative are differential, the rest algebraic:
\begin{equation}
  \text{differential: }\;
  \{\mathbf y_{\mathrm{cell}}\}\cup\{\mathbf y_{\mathrm{pc}}\},
  \label{eq:ida_diff}
\end{equation}
\begin{equation}
  \text{algebraic: }\;
  \{\mathbf y_{\mathrm{trace}}\}.
  \label{eq:ida_alg}
\end{equation}
These are communicated to IDA through the \texttt{differential\_components}
mask, which now flags two separate index ranges (the cell block and the trailing
$P_c$ block) as differential, with the algebraic trace block between them.
Appending the first-order capacitor ODE does not raise the index: the trace
block $\partial\mathbf R_{\mathrm{trace}}/\partial\mathbf y_{\mathrm{trace}}$
remains nonsingular, so~\eqref{eq:ida_residual} is a \emph{semi-explicit
index-1} DAE.

\paragraph{Newton system and Jacobian.}
At each step IDA supplies the BDF coefficient $\alpha$ and requests
\begin{equation}
  J_{\mathrm{IDA}}
  =\frac{\partial\mathbf F}{\partial\mathbf y}
   +\alpha\,\frac{\partial\mathbf F}{\partial\dot{\mathbf y}},
  \qquad
  \frac{\partial\mathbf F}{\partial\dot{\mathbf y}}
  =\mathrm{blockdiag}\bigl(M_K,\,\mathbf 0,\,\mathbf C_{\mathrm{wk}}\bigr),
  \quad
  \mathbf C_{\mathrm{wk}}=\mathrm{diag}(C_i).
  \label{eq:ida_J}
\end{equation}
This yields the $3\times3$ block structure
\begin{equation}
  J_{\mathrm{IDA}}
  =
  \begin{pmatrix}
  \alpha M_K + A & B & 0\\[4pt]
  C & D & E\\[4pt]
  0 & G & \alpha\,\mathbf C_{\mathrm{wk}}+\mathbf R_2^{-1}
  \end{pmatrix},
  \label{eq:ida_Jblocks}
\end{equation}
in which $A,B,C,D$ are the cell--cell, cell--trace, trace--cell and
trace--trace blocks already assembled for the network (the trace blocks contain
the per-face Jacobians, e.g.\ $D_{\mathrm{rcr}}$ of~\eqref{eq:Drcr} and
$C_{\mathrm{rcr}}$ of~\eqref{eq:Brcr}). The capacitor coupling reuses the
boundary-section derivatives directly:
\begin{equation}
  E \;\text{contains}\; E_{\mathrm{rcr}}
  =\frac{\partial\widehat{\mathbf B}_h^{\mathrm{rcr}}}{\partial P_c}
  =\begin{bmatrix}-1\\[2pt]0\end{bmatrix}
  \quad\text{from~\eqref{eq:Ercr}},
  \label{eq:ida_E}
\end{equation}
\begin{equation}
  G \;\text{contains}\;
  \frac{\partial R_{\mathrm{pc},i}}{\partial(\widehat A_i,\widehat U_i)}
  =\bigl(-\widehat U_i,\;-\widehat A_i\bigr)
  \quad\text{from~\eqref{eq:rcr_cap_jac}},
  \label{eq:ida_G}
\end{equation}
and the capacitor diagonal entry per outlet is
\begin{equation}
  \frac{\partial R_{\mathrm{pc},i}}{\partial P_{c,i}}
  +\alpha\,\frac{\partial R_{\mathrm{pc},i}}{\partial\dot P_{c,i}}
  =\frac{1}{R_{2,i}}+\alpha\,C_i ,
  \label{eq:ida_pcdiag}
\end{equation}
using~\eqref{eq:rcr_cap_jac} and~\eqref{eq:rcr_cap_mass}. The two mass terms
$\alpha M_K$ and $\alpha C_i$ are stamped in \texttt{assemble\_jacobian()}
(the latter as \texttt{jacobian\_matrix.add(pc\_dof, pc\_dof, alpha *
rcr\_map.at(bid).C)}); the remaining entries come from the residual Jacobian
assembly. Here the capacitance $C_i$ is distinct from the trace--cell block
$C$.

\paragraph{Linear solve.}
Each Newton iteration solves
\begin{equation}
  J_{\mathrm{IDA}}\,\delta\mathbf y=-\mathbf F
  \label{eq:ida_linsolve}
\end{equation}
with the sparse direct solver UMFPACK. Relative to the explicit-update
formulation only the system size and sparsity change: a few extra rows/columns
for $\mathbf y_{\mathrm{pc}}$ and the off-diagonal $E,G$ couplings.

\paragraph{Consistent initialisation.}
At $t=0$ IDA requires a pair $(\mathbf y_0,\dot{\mathbf y}_0)$ satisfying
$\mathbf F(0,\mathbf y_0,\dot{\mathbf y}_0)=\mathbf 0$. The cell DOFs are seeded
at rest,
\begin{equation}
  A=A_d(x),\qquad U=0,
  \label{eq:ida_seed_cell}
\end{equation}
and each capacitor pressure at a prescribed level $P_{c,i}(0)$ (e.g.\ diastolic,
or $P_{\mathrm{out},i}$). The trace DOFs are then obtained from a dedicated
Newton solve of
\begin{equation}
  \mathbf R_{\mathrm{trace}}(0,\mathbf y_0)=\mathbf 0
  \label{eq:ida_traceinit}
\end{equation}
holding the cell and capacitor DOFs fixed (routine
\texttt{initialize\_trace\_unknowns()}). Operationally this is enforced by
replacing the cell and $P_c$ rows of the Newton matrix with identity rows (unit
diagonal, zero right-hand side), so their increments vanish and they remain at
the seeded values, while the nonsingular trace block is inverted for
$\mathbf y_{\mathrm{trace}}$. Stamping the unit diagonal on the $P_c$ rows is
essential: those rows are otherwise empty in the trace-only system and the
Newton matrix would be singular.

The differential rates follow from the residual. The cell rate is
\begin{equation}
  \dot{\mathbf y}_{\mathrm{cell}}
  =M_K^{-1}\,\mathbf R_{\mathrm{cell}}(0,\mathbf y_0),
  \label{eq:ida_ydotcell}
\end{equation}
and, since $U=0\Rightarrow Q_h=0$, the capacitor rate from
$\mathbf R_{\mathrm{pc}}=\mathbf 0$ is
\begin{equation}
  \dot P_{c,i}(0)
  =-\,\frac{P_{c,i}(0)-P_{\mathrm{out},i}}{R_{2,i}\,C_i},
  \label{eq:ida_ydotpc}
\end{equation}
with $\dot{\mathbf y}_{\mathrm{trace}}=\mathbf 0$. IDA's built-in
consistent-IC correction (\texttt{ic\_type = use\_y\_diff}) then refines the
pair $(\mathbf y_0,\dot{\mathbf y}_0)$ before the first step.

\section{Example: Global HDG Matrix Structure for a Bifurcating Network}

Consider a simple arterial network consisting of:
\begin{itemize}
\item one parent vessel $P$,
\item two daughter vessels $D_1$ and $D_2$,
\item one inflow boundary,
\item one bifurcation junction,
\item two RCR outlet boundaries.
\end{itemize}

\subsection{Global Newton System}

After assembling all transmission residuals, the Newton linearisation
at each time step yields the block system
\begin{equation}
\begin{pmatrix}
A_P    & 0      & 0      & B_P     & 0   \\
0      & A_{D_1}& 0      & B_{D_1} & 0   \\
0      & 0      & A_{D_2}& B_{D_2} & 0   \\
C_P    & C_{D_1}& C_{D_2}& D       & E   \\
0      & 0      & 0      & G       & H
\end{pmatrix}
\begin{pmatrix}
\delta W_P \\ \delta W_{D_1} \\ \delta W_{D_2} \\ \delta\widehat{W} \\ \delta P_c
\end{pmatrix}
=
-
\begin{pmatrix}
R_P \\ R_{D_1} \\ R_{D_2} \\ R_{\widehat{W}} \\ R_{P_c}
\end{pmatrix}.
\label{eq:block_hdg_system}
\end{equation}

The trace unknown vector, ordered by face type, is
\begin{equation}
\widehat{W}
=
\begin{pmatrix}
\widehat{W}_{\mathrm{in}}   \\
\widehat{W}_{P,1}           \\
\widehat{W}_{P,2}           \\
\widehat{W}_{J}             \\
\widehat{W}_{D_1,1}         \\
\widehat{W}_{D_1,2}         \\
\widehat{W}_{\mathrm{out},1}\\
\widehat{W}_{D_2,1}         \\
\widehat{W}_{D_2,2}         \\
\widehat{W}_{\mathrm{out},2}
\end{pmatrix},
\qquad
P_c
=
\begin{pmatrix}
P_{c,1}\\
P_{c,2}
\end{pmatrix},
\end{equation}
where each $\widehat{W}$ entry is a 2-vector $(\widehat{A}, \widehat{U})^\top$
at the corresponding face, and $P_{c,1},P_{c,2}$ are the capacitor pressures of
the two RCR outlets, where
\begin{align}
A_\alpha
&= \frac{\partial R_\alpha}{\partial W_\alpha}
& &\longrightarrow \text{element residual w.r.t. element unknowns}
\\[6pt]
B_\alpha
&= \frac{\partial R_\alpha}{\partial \widehat{W}}
& &\longrightarrow \text{element residual w.r.t. trace unknowns}
\\[6pt]
C_\alpha
&= \frac{\partial R_{\widehat{W}}}{\partial W_\alpha}
& &\longrightarrow \text{transmission residual w.r.t. element unknowns}
\\[6pt]
D
&= \frac{\partial R_{\widehat{W}}}{\partial \widehat{W}}
& &\longrightarrow \text{transmission residual w.r.t. trace unknowns}
\\[6pt]
E
&= \frac{\partial R_{\widehat{W}}}{\partial P_c}
& &\longrightarrow \text{transmission residual w.r.t. capacitor pressures}
\\[6pt]
G
&= \frac{\partial R_{P_c}}{\partial \widehat{W}}
& &\longrightarrow \text{capacitor residual w.r.t. trace unknowns}
\\[6pt]
H
&= \frac{\partial R_{P_c}}{\partial P_c}
& &\longrightarrow \text{capacitor residual w.r.t. capacitor pressures}
\end{align}

\subsection{Element Residual Vectors}

The vessel residual vectors $R_\alpha$ appearing in the right-hand
side of~\eqref{eq:block_hdg_system} are assembled from the local
weak form residual on each element $K \in \alpha$:

\begin{equation}
R_\alpha
=
\begin{pmatrix}
R_{K_1} \\ R_{K_2} \\ R_{K_3}
\end{pmatrix},
\end{equation}

where for each element $K$ the local residual is:
\begin{equation}
R_K(\mathbf{w}_h, \widehat{\mathbf{w}}_h)
=
\underbrace{
\left(\frac{\mathbf{w}_h - \mathbf{w}_h^n}{\Delta t},
\mathbf{v}\right)_K
}_{\text{time derivative}}
-
\underbrace{
\left(\mathbf{F}(\mathbf{w}_h),
\bm{b}\cdot\nabla\mathbf{v}\right)_K
}_{\text{volume flux}}
+
\underbrace{
\left\langle
\widehat{\mathbf{F}}_h \cdot \bm{n},
\mathbf{v}
\right\rangle_{\partial K}
}_{\text{face flux via trace}}
-
\underbrace{
\left(\mathbf{S}(\mathbf{w}_h),
\mathbf{v}\right)_K
}_{\text{source}},
\label{eq:RK}
\end{equation}

with the numerical flux $\widehat{\mathbf{F}}_h$ defined
in~\eqref{eq:hdg_local}. The face flux term couples $R_K$
to the trace unknowns $\widehat{\mathbf{w}}_h$, which is
why the $B_\alpha$ blocks appear in the element rows of
the global system. The element residuals do not depend on
$P_c$, so the element rows of~\eqref{eq:block_hdg_system} have
a zero column under $\delta P_c$.

Similarly, $R_{\widehat{W}}$ is the vector of transmission
residuals assembled from all faces:
\begin{equation}
R_{\widehat{W}}
=
\begin{pmatrix}
\mathcal{F}_{\mathrm{in}}(\mathbf{w}_h, \widehat{W}_{\mathrm{in}})              \\
\mathcal{F}_{f}(\mathbf{w}_h, \widehat{W}_{P,1})                                \\
\mathcal{F}_{f}(\mathbf{w}_h, \widehat{W}_{P,2})                                \\
\mathcal{F}_{J}(\widehat{W}_{J})                                                 \\
\mathcal{F}_{f}(\mathbf{w}_h, \widehat{W}_{D_1,1})                              \\
\mathcal{F}_{f}(\mathbf{w}_h, \widehat{W}_{D_1,2})                              \\
\mathcal{F}_{\mathrm{rcr}}(\mathbf{w}_h, \widehat{W}_{\mathrm{out},1}, P_{c,1}) \\
\mathcal{F}_{f}(\mathbf{w}_h, \widehat{W}_{D_2,1})                              \\
\mathcal{F}_{f}(\mathbf{w}_h, \widehat{W}_{D_2,2})                              \\
\mathcal{F}_{\mathrm{rcr}}(\mathbf{w}_h, \widehat{W}_{\mathrm{out},2}, P_{c,2})
\end{pmatrix},
\label{eq:R_What}
\end{equation}

each entry being the transmission residual, with the two RCR rows also
carrying the corresponding capacitor pressure $P_{c,i}$.

The capacitor residual vector collects one Windkessel
equation~\eqref{eq:rcr_capacitor} per RCR outlet:
\begin{equation}
R_{P_c}
=
\begin{pmatrix}
R_{\mathrm{pc},1}\\
R_{\mathrm{pc},2}
\end{pmatrix},
\qquad
R_{\mathrm{pc},i}
=
C_i\,\frac{P_{c,i}-P_{c,i}^{\,n}}{\Delta t}
-\widehat{A}_i\widehat{U}_i
+\frac{P_{c,i}-P_{\mathrm{out},i}}{R_{2,i}} .
\label{eq:R_Pc}
\end{equation}

\subsection{Structure of the Coupling Blocks}

Each vessel-level coupling matrix $C_\alpha$ is assembled from
the individual face transmission Jacobians $C_{\mathrm{in}}$,
$C_f$, $C_J$, $C_{\mathrm{rcr}}$ defined in
equations~\eqref{eq:inflow_jacobian}--\eqref{eq:Cjunc}.
Each row block corresponds to the transmission condition at one
face of the vessel, and acts on the element DOFs $W_\alpha$
in the column direction, here $C_J =0$

\begin{equation}
C_P
=
\renewcommand{\arraystretch}{1.6}
\left(
\begin{array}{c}
C_{\mathrm{in}}  \\
\hline
C_{f}            \\
\hline
C_{f}            \\
\hline
C_{J}
\end{array}
\right)
\begin{array}{l}
\leftarrow \text{inflow face } \widehat{W}_{\mathrm{in}}     \\
\leftarrow \text{interior face } \widehat{W}_{P,1}           \\
\leftarrow \text{interior face } \widehat{W}_{P,2}           \\
\leftarrow \text{junction face } \widehat{W}_{J}
\end{array}
\label{eq:CP}
\end{equation}

\begin{equation}
C_{D_1}
=
\renewcommand{\arraystretch}{1.6}
\left(
\begin{array}{c}
C_{J}            \\
\hline
C_{f}            \\
\hline
C_{f}            \\
\hline
C_{\mathrm{rcr}}
\end{array}
\right)
\begin{array}{l}
\leftarrow \text{junction face } \widehat{W}_{J}             \\
\leftarrow \text{interior face } \widehat{W}_{D_1,1}         \\
\leftarrow \text{interior face } \widehat{W}_{D_1,2}         \\
\leftarrow \text{RCR outlet face } \widehat{W}_{\mathrm{out},1}
\end{array}
\label{eq:CD1}
\end{equation}

\begin{equation}
C_{D_2}
=
\renewcommand{\arraystretch}{1.6}
\left(
\begin{array}{c}
C_{J}            \\
\hline
C_{f}            \\
\hline
C_{f}            \\
\hline
C_{\mathrm{rcr}}
\end{array}
\right)
\begin{array}{l}
\leftarrow \text{junction face } \widehat{W}_{J}             \\
\leftarrow \text{interior face } \widehat{W}_{D_2,1}         \\
\leftarrow \text{interior face } \widehat{W}_{D_2,2}         \\
\leftarrow \text{RCR outlet face } \widehat{W}_{\mathrm{out},2}
\end{array}
\label{eq:CD2}
\end{equation}

\subsection{Structure of the Capacitor Coupling Blocks $E$ and $G$}

The block $E=\partial R_{\widehat{W}}/\partial P_c$ has nonzero rows only at
the two RCR outlet faces, each given by $E_{\mathrm{rcr}}$
of~\eqref{eq:Ercr}; all other trace rows are zero:
\begin{equation}
E
=
\renewcommand{\arraystretch}{1.6}
\left(
\begin{array}{cc}
0              & 0              \\
\vdots         & \vdots         \\
E_{\mathrm{rcr}} & 0            \\
\vdots         & \vdots         \\
0              & E_{\mathrm{rcr}}
\end{array}
\right)
\begin{array}{l}
\\
\leftarrow \text{RCR outlet face } \widehat{W}_{\mathrm{out},1}\\
\\
\leftarrow \text{RCR outlet face } \widehat{W}_{\mathrm{out},2}
\end{array},
\qquad
E_{\mathrm{rcr}}
=
\begin{bmatrix}-1\\[2pt]0\end{bmatrix}.
\label{eq:E_block}
\end{equation}

The block $G=\partial R_{P_c}/\partial\widehat{W}$ has nonzero columns only at
the two RCR outlet trace faces:
\begin{equation}
G
=
\renewcommand{\arraystretch}{1.6}
\begin{pmatrix}
\cdots & G_1 & \cdots & 0   & \cdots \\
\cdots & 0   & \cdots & G_2 & \cdots
\end{pmatrix},
\qquad
G_i
=
\frac{\partial R_{\mathrm{pc},i}}{\partial(\widehat{A}_i,\widehat{U}_i)}
=
\bigl(-\widehat{U}_i,\;-\widehat{A}_i\bigr),
\label{eq:G_block}
\end{equation}
the column of $G_i$ aligning with the $\widehat{W}_{\mathrm{out},i}$ entry of
$\widehat{W}$, from~\eqref{eq:rcr_cap_jac}.

\subsection{Structure of the Trace Block $D$}

The trace block $D$ is block-diagonal, with each diagonal entry
being the Jacobian $D_F = \partial\mathcal{F}_F/\partial
\widehat{\mathbf{w}}_h$ at the corresponding face:

\begin{equation}
D
=
\mathrm{blockdiag}
\bigl(
D_{\mathrm{in}},\;
D_{f},\;
D_{f},\;
D_{J},\;
D_{f},\;
D_{f},\;
D_{\mathrm{rcr}},\;
D_{f},\;
D_{f},\;
D_{\mathrm{rcr}}
\bigr),
\label{eq:D_blockdiag}
\end{equation}
with the $2\times2$ values:
\begin{equation}
\renewcommand{\arraystretch}{1.8}
\begin{array}{ll}
D_{\mathrm{in}}
=
\begin{bmatrix}
\widehat{U}_h & \widehat{A}_h \\
4\,dc/d\widehat{A}_h & 1
\end{bmatrix},
&
\text{(inflow face)}
\\[10pt]
D_{f}
=
-2\alpha_{\mathrm{LF}}\,\mathbf{I},
&
\text{(interior face)}
\\[10pt]
D_{\mathrm{rcr}}
=
\begin{bmatrix}
dP/d\widehat{A}_h - R_1\widehat{U}_h
& -R_1\widehat{A}_h \\
4\,dc/d\widehat{A}_h & 1
\end{bmatrix},
&
\text{(RCR outlet face)}
\end{array}
\label{eq:D_blocks_summary}
\end{equation}
and $D_J$ is the $3\times6$ block from equation~\eqref{eq:DJ},
spanning the three junction trace unknowns
$(\widehat{W}_{J,P},\, \widehat{W}_{J,D_1},\, \widehat{W}_{J,D_2})$.
With $P_c$ a global unknown, dependence is now carried by the off-diagonal block
$E$ through $E_{\mathrm{rcr}}$.

\subsection{Structure of the Capacitor Block $H$}

The capacitor block $H=\partial R_{P_c}/\partial P_c$ is diagonal, one scalar
entry per RCR outlet,
\begin{equation}
H
=
\mathrm{diag}\!\left(
\alpha C_1+\frac{1}{R_{2,1}},\;
\alpha C_2+\frac{1}{R_{2,2}}
\right),
\label{eq:H_block}
\end{equation}
from~\eqref{eq:rcr_cap_jac} and~\eqref{eq:rcr_cap_mass} (under IDA,
$C_i/\Delta t \to \alpha\,C_i$ with the BDF coefficient $\alpha$).
\bibliographystyle{plain}
\bibliography{metric_flow}
\end{document}